\documentclass[11pt,twocolumn,letterpaper]{article}
\usepackage[spanish,es-nodecimaldot]{babel}
\usepackage[utf8]{inputenc}
\usepackage[T1]{fontenc}
\usepackage{float}

%% Sets page size and margins
\usepackage[a4paper,top=2.5cm,bottom=2.5cm,left=2cm,right=2cm]{geometry}

\usepackage{amsmath}
\usepackage{amssymb}
\usepackage{amsthm}
\usepackage{mathtools}
\usepackage{listings}
\usepackage{enumitem}[shortlabels]
\usepackage[colorinlistoftodos]{todonotes}
\usepackage[colorlinks=true, allcolors=blue]{hyperref}
\usepackage{url}

\usepackage{lipsum}

\usepackage{tcolorbox}

%Author affil
\usepackage{authblk}

%% Title
\title{
		\vspace{-0.7in} 	
		\usefont{OT1}{bch}{b}{n}
		
		\begin{minipage}{3cm}
        \vspace{-0.5in} 	
    	\begin{center}
    		\includegraphics[height=3.2cm]{logo_unam.png}
    	\end{center}
    \end{minipage}\hfill
    \begin{minipage}{10.7cm}
    	\begin{center}

\normalfont \normalsize \textsc{UNIVERSIDAD NACIONAL AUTÓNOMA DE MÉXICO \\ FACULTAD DE CIENCIAS \\ LABORATORIO} \\
		\huge Número de práctica \\ Nombre de Práctica\\
        
        \normalsize \rmfamily \mdseries
        \vspace{0.1in}
        Autor$^1$\\

        \vspace{0.1in}
        $^1$Facultad de Ciencias, UNAM\\
        \date{\date{\vspace{0.1in}\today}}
    	\end{center}
     
    \end{minipage}\hfill
    \begin{minipage}{3.2cm}
    \vspace{-0.5in} 
    	\begin{center}
    		\includegraphics[height=3.2cm]{logo_fc.png}
    	\end{center}
    \end{minipage}

\author{}

\vspace{0.1in}
\begin{tcolorbox}[colframe=white, arc=8pt]
\normalsize \rmfamily \mdseries
\centering
\textbf{Resumen}\\

% Puedes colocar el texto del resumen en esta parte

\lipsum[2]

\textit{Palabras clave:} \textbf{Lorem, Ipsum, eusimod, sedes }
\end{tcolorbox}
\vspace{-0.7in} 
}




\hypersetup{
    colorlinks=true,
    linkcolor=blue,
    filecolor=magenta,      
    urlcolor=cyan
    }


%------Vectores Unitarios-------%
\newcommand{\uveci}{\hat{\textbf{i}}} 
\newcommand{\uvecj}{\hat{\textbf{j}}}
\newcommand{\uveck}{\hat{\textbf{k}}}

\newcommand{\grade}{^{\circ}}
 \setlength {\marginparwidth }{2cm}
 
\begin{document}

\maketitle


\section{Objetivos}
\begin{itemize}
    \item Nulla malesuada porttitor diam. Donec felis
erat, congue non, volutpat at, tincidunt tristique, libero. 
    \item  Nulla malesuada porttitor diam. Donec felis
erat, congue non, volutpat at, tincidunt tristique, libero. 

    \item  Nulla malesuada porttitor diam. Donec felis
erat, congue non, volutpat at, tincidunt tristique, libero. 
\end{itemize}

\section{Introducción}

% Introducción: porque hacer el experimento, para medir que, etc.

\subsection{Subtema 1}

\lipsum[2] \cite{sears}.

\begin{figure}[h!]
            \centering
            \begin{tikzpicture}[scale = 0.9]
                \draw[thick,fill = gray!40, opacity = 0.1] (-4,-3) rectangle (4,3);

                \draw[thick, fill = gray!40, opacity = 0.5, dashed] (0,0) circle (2.5);

                

                \draw[thick, fill = gray!40, opacity = 0.7, dashed] (1,1) circle (1.07);

                \draw[dashed] (0,0) -- (1.77,1.77);

                \draw[dashed] (0,0) -- (-2.5,0) node[midway, above]{$r$};

                \draw[double] (1,1) -- (1.77,1.77) node[midway, above left]{$\epsilon$};

                \draw[|-|] (-0.2,0.2) -- (0.8,1.2);

                \node[above,rotate=45,scale=0.6] at (0.5,0.8) {$d(x,c)$};

                \node at (2.5,2.5) {$M$};
                
                \fill[fill=black] (0,0) circle (1.3pt) node[below left]{$c$};

                \fill[fill=black] (1,1) circle (1.3pt) node[below right]{$x$};

                %\node[above left] at (0.5,0.5) {$r$};

                \end{tikzpicture}
            \caption{Descripción y dibujos con tikz}
        \end{figure}

\subsection{Subtema 2}
Ecuaciones \cite{Neil-2017}.

\begin{equation}
    I(p) = \left|\int \frac{exp(ikr(x))}{r(x)}\right|
\end{equation}
Ecuación:
\begin{equation}
    \frac{k\alpha sin(\theta)}{2}=\pm n \pi \quad n \in \mathbb{N}
\end{equation}




\section{Materiales}
\begin{itemize}
            \item Material 1
            \item Material 2
            \item Material 3
            \item Material 4


\end{itemize}





\section{Disposición y desarrollo experimental}
\subsection{Descripción}
Se pueden colocar imágenes

\begin{figure}[h!]
    \centering
    \includegraphics[width=0.8\linewidth]{img/disp_9.png}
    \caption{Descripción}
\end{figure}

\lipsum[3]\cite{griffiths-2017}

\section{Análisis de resultados}

\subsection{Subtema}

\lipsum[4]
\begin{table}[h!]
\centering
\begin{tabular}{|c|c|c|cc|}
\hline
{[}\#{]} & D {[}mm{]}     & $\Delta D$ {[}mm{]}        \\ \hline
1        & 780 $\pm$ 5    & 5.6 $\pm$ 0.327  \\ \hline
2        & 790 $\pm$ 5    & 4.82 $\pm$ 0.327 \\ \hline
3        & 780 $\pm$ 5    & 5.28 $\pm$ 0.327 \\ \hline
4        & 780 $\pm$ 5    & 4.6 $\pm$ 0.327  \\ \hline
5        & 780 $\pm$ 5    & 4.98 $\pm$ 0.327 \\ \hline
6        & 790 $\pm$ 5    & 4.62 $\pm$ 0.327 \\ \hline
7        & 790 $\pm$ 5    & 4.64 $\pm$ 0.327 \\ \hline
8        & 790 $\pm$ 5    & 4.8 $\pm$ 0.327  \\ \hline
9        & 790 $\pm$ 5    & 4.6 $\pm$ 0.327  \\ \hline
10       & 780 $\pm$ 5    & 4.56 $\pm$ 0.327 \\ \hline
Promedio & 785 $\pm$ 1.58 & 4.85 $\pm$ 0.103 \\ \hline
\end{tabular}
\caption{Descripción de tabla}
\end{table}

\lipsum[6]\cite{halliday-2002}

\begin{figure}[h!]
    \centering
    \includegraphics[width=0.9\linewidth]{img/res_1.png}
    \caption{Graficas}
\end{figure}




\section{Discusión y conclusiones}
\lipsum[5]


\bibliographystyle{unsrt}

%Referencias en formato bibtex
\bibliography{references.bib}

\clearpage


\end{document}